\section{Background}
Contracts are part of almost everything people and businesses do, but reading them takes a lot of time. A typical contract in our dataset is about 33{,}000 characters long (roughly 5{,}000 words), and the longest is more than 338{,}000 characters. Somewhere in all that text are the parts that really matter: who owns new intellectual property, whether there is a limit on how much you can owe, whether the contract renews on its own, and whether you are stopped from working with competitors after it ends. Having a lawyer read a contract takes hours and costs money that many people and small businesses cannot afford. So people often sign contracts they do not fully understand and take on risks they never noticed.

The Contract Understanding Atticus Dataset (CUAD) \cite{hendrycks2021cuad} was created to help automate this kind of review. It has 510 real business contracts in which legal experts marked every clause that belongs to one of 41 important categories. This marking work is said to be worth about two million US dollars of lawyers' time. However, CUAD only treats the problem as finding text: given a category, find the clause. Once it finds a clause such as ``Uncapped Liability'', it stops. It does not warn the reader what the clause means for them, it does not explain it in simple words, and it does not check how clauses affect each other.

\subsection{System at a Glance}
Chapter~4 explains the system in detail. Figure~\ref{fig:overview} shows the whole system in one picture. The system answers three questions about a contract, in order: \emph{what is in it}, \emph{where is it}, and \emph{what does it mean}. CUAD was built to test only the first two. The third step, shaded in the figure, is what this project adds.

\begin{figure}[ht]
\centering
\begin{adjustbox}{max width=\textwidth}
\begin{tikzpicture}[
  node distance=6mm,
  box/.style={draw, rounded corners=2pt, minimum height=19mm, text width=25mm,
              align=center, font=\small, fill=blue!6},
  new/.style={box, fill=orange!15},
  arr/.style={-{Stealth[length=2.5mm]}, thick}
]
\node[box] (a) {\textbf{Contract}\\[2pt]{\scriptsize raw text}\\{\scriptsize median 33k chars}};
\node[box, right=of a] (b) {\textbf{What is in it?}\\[2pt]{\scriptsize presence over}\\{\scriptsize 41 categories}};
\node[box, right=of b] (c) {\textbf{Where is it?}\\[2pt]{\scriptsize exact clause}\\{\scriptsize text extracted}};
\node[new, right=of c] (d) {\textbf{What does it mean?}\\[2pt]{\scriptsize risk level +}\\{\scriptsize plain English}};
\draw[arr] (a) -- (b);
\draw[arr] (b) -- (c);
\draw[arr] (c) -- (d);

\draw[decorate, decoration={brace, mirror, amplitude=4pt}]
  ($(b.south west)+(0,-1.5mm)$) -- ($(c.south east)+(0,-1.5mm)$)
  node[midway, below=4pt, font=\scriptsize\itshape] {CUAD's original task};
\draw[decorate, decoration={brace, mirror, amplitude=4pt}]
  ($(d.south west)+(0,-1.5mm)$) -- ($(d.south east)+(0,-1.5mm)$)
  node[midway, below=4pt, font=\scriptsize\itshape] {added here};
\end{tikzpicture}
\end{adjustbox}
\caption{Three questions the system answers}
\label{fig:overview}
\end{figure}

\section{Rationale of the Study or Motivation}
Our project tries to close the gap between finding a clause and understanding it. On top of finding clauses, we add two things. First, a simple risk table that gives every clause \emph{category} a High, Medium or Low level, judged from the side of the weaker party. Second, a text model that is meant to turn each clause into one plain-English sentence. We built the clause-finding models and the risk table into \emph{Clausify}, a prototype web app: a user pastes or uploads a contract and gets back a list of clauses sorted by risk. Both added parts worked less well than we hoped. The risk table has not been checked by a lawyer (Section~\ref{sec:taxvalid}), and the text model learned to repeat a fixed sentence for each category instead of describing the actual clause (Section~\ref{sec:summ}). Chapters~5 and~6 report these results honestly.

So the reason for this study is practical as well as academic. Earlier work on CUAD answers the question ``where is the clause?'' for a lawyer who already knows what to look for. The people who need help most (individuals, freelancers and small businesses who sign contracts without a lawyer) also need to know which clauses matter most to them and what each one means in everyday language. A system that gives those answers, and is honest about how often it is wrong, could make contract review possible for people who cannot pay for it today. Just as important, we wanted to measure how well such a system really works instead of assuming it. A contract-review tool that looks confident but quietly misses clauses is worse than having no tool at all.

\section{Problem Statement}
In short, the project works on four connected problems:
\begin{enumerate}[itemsep=2pt]
  \item \textbf{Finding which clauses are present.} For each of the 41 CUAD categories, decide whether a contract has at least one clause of that type. The main difficulty is length: a typical contract is about 16$\times$ longer than the amount of text a standard transformer model can read at once.
  \item \textbf{Finding the exact clause text.} For each category that is present, find the exact words of the clause in the contract. We treat this as a question-answering task, in the same style as SQuAD \cite{rajpurkar2016squad}.
  \item \textbf{Explaining in plain language.} Rewrite each clause as one sentence that a non-lawyer can understand.
  \item \textbf{Showing the risk.} Give every clause found a risk level and a one-line reason, so the result reads like a ranked list of warnings instead of a list of quotes. We did not test whether this actually helps users (Section~\ref{sec:threats}).
\end{enumerate}

\subsection{Research Questions}
\label{sec:rqs}
Our objectives turn into five research questions, which the rest of this report answers. We give the answers here so readers can check each claim against the evidence instead of searching for it.

\begin{description}[itemsep=4pt, leftmargin=1.6cm]
  \item[RQ1] \emph{Can a transformer that reads a contract in small windows beat a simple word-based model that reads the whole contract, at finding which clauses are present?}\\
  \textbf{No.} The simple model wins by $+0.082$ micro-F1. The gap stays even when the transformer is trained on all its training data (Section~\ref{sec:budget}) and with three different random seeds (Section~\ref{sec:seeds}). Answered in Chapter~5.
  \item[RQ2] \emph{How well can the system find the exact clause text in long contracts?}\\
  Token-F1 is 0.764 (95\% CI $[0.743, 0.782]$) \emph{when the model is given a window that contains the clause}, but only 0.383 when the window is chosen by the first model (Sections~\ref{sec:span} and~\ref{sec:e2e}).
  \item[RQ3] \emph{Does the summarizer really describe each specific clause in plain English?}\\
  \textbf{No.} For every input it copies one of 41 sentences we wrote in advance. Against clause-specific reference sentences, it scores no better than just printing that fixed sentence without any model (Section~\ref{sec:summ}).
  \item[RQ4] \emph{How do errors add up across the whole system?}\\
  They multiply, and it is worse than the separate scores suggest: only 22.0\% of real clauses get through all three steps, compared with the $\approx$65\% that the separate scores would predict (Section~\ref{sec:e2e}).
  \item[RQ5] \emph{Do the system's mistakes avoid the clauses that would hurt a signer most?}\\
  \textbf{No, the opposite.} The setup we report as our main result misses 36.9\% of High-risk clauses, while the model we call weaker misses only 15.9\% (Section~\ref{sec:riskeval}). The overall score and the harm to the user point in different directions.
\end{description}

Three of these five answers are negative. We report them as real results of the project, not as failures. Each one is something a reader would have got wrong by looking only at the separate model scores.

\section{Objective}
\begin{itemize}[itemsep=2pt]
  \item Build and fairly test a complete three-model CUAD system on an ordinary laptop, using an 80/20 split by contract so that every reported number comes from contracts the models have never seen.
  \item Check whether a transformer really beats a strong simple model at finding which clauses are present, instead of assuming it does.
  \item Measure how well the exact clause text is found, using token-level F1 and exact match on the test set.
  \item Test the summarizer against two different sets of reference sentences, because one automatic score turned out to measure the reference sentences rather than the task.
  \item Measure how many clauses survive the whole system from start to finish, not only the score of each step.
  \item Put the trained models into a working demo app that shows clauses sorted by risk.
\end{itemize}

\section{Methodology in Brief}
We tested every design before building on it, using one fixed 80/20 split of CUAD by contract (408 training contracts and 102 test contracts). After exploring the data (Section~\ref{sec:eda}), we first built a simple word-based model (TF--IDF with logistic regression) that reads the whole contract, \emph{before} training any transformer. We also measured the highest recall a ``search first'' design could ever reach (Section~\ref{sec:ceiling}). Based on that, we trained a DistilBERT model that scores the contract in small windows and takes the highest window score as the answer for the whole contract (multiple-instance learning). We combined it with the simple model using an AND rule, which reports a clause only when both models agree. A DistilBERT question-answering model finds the exact clause text, a fine-tuned FLAN-T5-small model writes a one-sentence plain-English explanation, and a hand-made table gives each of the 41 categories a High, Medium or Low risk level.

Every model was tested once on the 102 test contracts. We report each score together with confidence intervals from a bootstrap over contracts, and we checked how the results change with the amount of training data, the random seed, the way window scores are combined, the decision threshold and the train/test split. We also tested the full system from start to finish. Finally, we put the trained models into \emph{Clausify}, a Streamlit web app that gives a risk-sorted clause report for any contract. Chapter~4 describes the method in full.

\section{Scopes and Challenges}
\textbf{Scope.} The project covers English business contracts like those in CUAD: 510 contracts from US SEC filings, marked for 41 clause categories. Within that scope it handles four tasks (finding which clauses are present, finding their exact text, explaining them in plain language and showing their risk level) and delivers a working prototype app. The following are outside the scope on purpose: legal advice of any kind; contracts in other languages or countries (including Bangladeshi contracts, which inspired the work but were not tested); judging the risk of a clause's exact wording; finding conflicts between clauses; and a formal user study of the app. All models were sized so they can be trained and run on one ordinary laptop.

\textbf{Challenges.} Six challenges shaped the design:
\begin{enumerate}[itemsep=2pt]
  \item \textbf{Very long documents.} A typical contract is about 33{,}000 characters. That is roughly 16$\times$ the 512-token limit of a standard transformer (about 2{,}000 characters), and the longest contract is more than 338{,}000 characters, about 170$\times$.
  \item \textbf{Very uneven categories.} Some categories appear in 100\% of contracts and some in only 2.5\%; a few have only a handful of examples in the whole dataset.
  \item \textbf{A small, fixed dataset.} The 510 contracts are all there is, and expert marking is too costly to add more, so we had to be careful about over-fitting and about test data leaking into training.
  \item \textbf{Limited hardware.} Training ran on an Apple-Silicon laptop without a CUDA GPU. Our first-choice model (DeBERTa-v3) was unstable on the MPS backend, and GPU memory was not freed between training runs.
  \item \textbf{No plain-language answers to learn from.} CUAD has no plain-English versions of clauses, so we had to write our own reference sentences, and this limited what the summarizer could learn.
  \item \textbf{No agreed answer for risk either.} Contract risk depends on which side you are on, the situation and the country, and there is no expert-checked risk model for CUAD's 41 categories.
\end{enumerate}

\section{Report Organization}
The rest of this report is organized as follows. Chapter~2 reviews background and earlier work. Chapter~3 sets out the system's requirements, its effect on society, the environment, ethics and cost, and the project plan. Chapter~4 explains the method: how we designed the system, which other designs we considered, how the data was prepared, and how the models and the Clausify app were built. Chapter~5 analyses the results on the test set, including the statistics and the changes we made because of them. Chapter~6 sums up the findings, contributions and limits, and suggests future work. The appendices contain the full 41-category risk table, a checklist for repeating the work and the layout of the project files.

